% !TEX program = pdflatex
% ════════════════════════════════════════════════════════════════════════════
%  UrbanFlow CDMX / VialAI — Estrategia de negocio
%  Diplomado en Ciencia de Datos · FES Acatlán, UNAM · Gen 30 · 2026
%  Autor: Carlos Armando López Encino
%  Profesor: Mtro. Fernando Barranco Rodríguez
% ════════════════════════════════════════════════════════════════════════════

\documentclass[11pt,a4paper]{article}

% ───── Codificación y lenguaje ─────────────────────────────────
\usepackage[utf8]{inputenc}
\usepackage[T1]{fontenc}
\usepackage[spanish,es-tabla,es-noindentfirst]{babel}
\usepackage{lmodern}
\usepackage{microtype}

% ───── Matemáticas ─────────────────────────────────────────────
\usepackage{amsmath,amssymb}
\usepackage{siunitx}
\sisetup{output-decimal-marker={.},group-separator={\,}}

% ───── Geometría y espaciado ───────────────────────────────────
\usepackage[a4paper,margin=2.5cm,headheight=15pt]{geometry}
\usepackage{setspace}
\setstretch{1.15}
\usepackage{parskip}

% ───── Color e identidad VialAI ────────────────────────────────
\usepackage{xcolor}
\definecolor{vialNavyDeep}{HTML}{0B1F3A}
\definecolor{vialNavy}{HTML}{0F2E4D}
\definecolor{vialTeal}{HTML}{2EBFA6}
\definecolor{vialAqua}{HTML}{3BA3BD}
\definecolor{vialGreen}{HTML}{A3D88F}
\definecolor{vialAmber}{HTML}{F5B836}
\definecolor{vialText}{HTML}{1F2937}
\definecolor{vialMuted}{HTML}{4A5B6D}
\definecolor{vialBgSoft}{HTML}{F3F6FA}

% ───── Gráficos ────────────────────────────────────────────────
\usepackage{graphicx}
\graphicspath{{./}{imagenes_12/}{imagenes_12/imagenes_png/}}
\usepackage{float}

% ───── Tablas ──────────────────────────────────────────────────
\usepackage{booktabs}
\usepackage{array}
\usepackage{tabularx}
\usepackage{multirow}
\usepackage{colortbl}

% ───── Tipografía de secciones ─────────────────────────────────
\usepackage{titlesec}
\titleformat{\section}
  {\color{vialNavy}\normalfont\Large\bfseries}{\thesection.}{0.5em}{}
\titleformat{\subsection}
  {\color{vialNavy}\normalfont\large\bfseries}{\thesubsection}{0.5em}{}
\titleformat{\subsubsection}
  {\color{vialTeal}\normalfont\normalsize\bfseries}{\thesubsubsection}{0.5em}{}

% ───── Encabezados y pies ──────────────────────────────────────
\usepackage{fancyhdr}
\pagestyle{fancy}
\fancyhf{}
\fancyhead[L]{\color{vialMuted}\small VialAI --- Estrategia de negocio}
\fancyhead[R]{\color{vialMuted}\small\thepage}
\renewcommand{\headrulewidth}{0.4pt}
\renewcommand{\headrule}{\color{vialTeal}\hrule width\headwidth}

% ───── Cajas de callout ────────────────────────────────────────
\usepackage[most]{tcolorbox}

\newtcolorbox{tipbox}[1][]{%
  enhanced, breakable,
  colback=vialBgSoft, colframe=vialTeal, coltitle=vialNavy,
  fonttitle=\bfseries, title=#1,
  boxrule=0.8pt, arc=3mm, left=4mm, right=4mm, top=2mm, bottom=2mm,
}

\newtcolorbox{keybox}[1][]{%
  enhanced, breakable,
  colback=vialNavy!5, colframe=vialNavy, coltitle=white,
  fonttitle=\bfseries, title=#1,
  boxrule=1pt, arc=3mm, left=4mm, right=4mm, top=2mm, bottom=2mm,
}

\newtcolorbox{casebox}[1][]{%
  enhanced, breakable,
  colback=vialAmber!10, colframe=vialAmber!80!black, coltitle=vialNavy,
  fonttitle=\bfseries, title=#1,
  boxrule=0.8pt, arc=3mm, left=4mm, right=4mm, top=2mm, bottom=2mm,
}

% ───── Hipervínculos ───────────────────────────────────────────
\usepackage[colorlinks=true,
            linkcolor=vialNavy,
            citecolor=vialTeal,
            urlcolor=vialAqua]{hyperref}

% ───── Comandos propios ────────────────────────────────────────
\newcommand{\VialAI}{\textbf{\textcolor{vialTeal}{Vial}\textcolor{vialAmber}{AI}}}
\newcommand{\ZMVM}{\textsc{zmvm}}
\newcommand{\CDMX}{\textsc{cdmx}}

% ──────────────────────────────────────────────────────────────
\begin{document}

% ═══════════════════════ PORTADA ═══════════════════════════════
\begin{titlepage}
\centering
\vspace*{1.5cm}

{\color{vialMuted}\large Diplomado en Ciencia de Datos \textbullet\ Generación 30}\\[0.3em]
{\color{vialMuted}\small FES Acatlán \textbullet\ UNAM \textbullet\ 2026}\\[2cm]

{\color{vialTeal}\Huge\bfseries UrbanFlow \CDMX\ / \VialAI}\\[0.8em]
{\color{vialNavy}\LARGE\bfseries Estrategia de negocio}\\[0.5em]
{\color{vialMuted}\large\itshape Propuesta de valor y modelo de adopción\\
para el mercado de última milla en la \ZMVM}\\[2.5cm]

{\large\textbf{Carlos Armando López Encino}}\\
{\small\color{vialMuted}cale.delta@gmail.com}\\[0.6em]
{\small Profesor: Mtro. Fernando Barranco Rodríguez}\\[2cm]

{\color{vialMuted}\small Proyecto integrador \textbullet\ 18 de abril de 2026}\\[0.4em]
{\color{vialAqua}\small\url{https://github.com/caledelta/UrbanFlow_CDMX}}

\vfill
{\color{vialMuted}\footnotesize
Documento complementario al artículo técnico \texttt{proyecto\_ZMCDMX\_v3.pdf}.
Cubre el criterio ``Estrategia de negocio'' de la rúbrica de evaluación.}
\end{titlepage}

\newpage
\tableofcontents
\newpage

% ═══════════════════════ §1 EL PROBLEMA ════════════════════════
\section{El problema de negocio}

La \ZMVM\ ocupa el puesto~\textbf{7} en el TomTom Traffic Index 2024
entre las ciudades más congestionadas del mundo, con un conductor
promedio que pierde \num{119} horas al año atrapado en tráfico. Para
el automovilista casual, esa cifra es molesta. Para el conductor
profesional de última milla, es \emph{pérdida directa de ingresos}.

El mercado mexicano de entregas a domicilio alcanzó
\$\num{7.8}~mil~millones de dólares en 2024, con la \ZMVM\ concentrando
el \SI{38}{\percent} del volumen nacional. Aproximadamente
\num{180000} conductores de entrega operan diariamente en esta ciudad.
Su principal dolor: \textbf{no saber cuánto tardará realmente su
próximo viaje}. Las aplicaciones de navegación existentes entregan una
estimación puntual (\emph{``llegarás en 23 min''}), pero no comunican
la incertidumbre de esa cifra. Un conductor que acepta un pedido de
alta distancia sin saber si tardará 23 minutos o 47 minutos arriesga
perder bonos por SLA, quedar fuera de la ventana de entrega, o
encadenar mal sus rutas.

\VialAI\ resuelve exactamente ese vacío: en lugar de una estimación
puntual, entrega una \textbf{distribución completa de probabilidad}
---bandas $P_{10}$, $P_{50}$, $P_{90}$--- más un agente conversacional
que traduce esos números a una recomendación accionable en lenguaje
natural.

% ═══════════════════════ §2 PROPUESTA DUAL ═════════════════════
\section{Propuesta de valor: el sistema de acción dual}

\VialAI\ genera valor en dos direcciones simultáneas:

\begin{keybox}[Sistema de acción dual]
\textbf{Para el conductor profesional:}\\
Obtiene, antes de salir, una respuesta concreta a la pregunta
``¿vale la pena salir ahora, o conviene esperar 20 minutos?''.
El agente \VialAI\ razona sobre la distribución completa y entrega una
recomendación accionable en lenguaje natural, con la opción de
recibirla por voz para no dejar de conducir.\\[0.5em]
\textbf{Para el operador de flota:}\\
Obtiene un panel que muestra, para cada ruta activa, el Índice de
Confiabilidad $(IC = (P_{90} - P_{10}) / P_{50})$ en semáforo
verde/amarillo/rojo. Con ese panel puede decidir cuándo despachar,
qué rutas evitar, y cuándo aplicar perturbaciones conocidas
(manifestaciones, cierres del Metro) a su planificación.
\end{keybox}

\subsection{Para el conductor profesional}

El conductor interactúa con \VialAI\ de tres formas equivalentes:

\begin{enumerate}
\item \textbf{Manual:} selecciona origen y destino desde menús
desplegables, configura la hora de salida y pulsa
\emph{``Predecir trayecto''}.
\item \textbf{Chat en lenguaje natural:} escribe
\emph{``¿Cuánto tardo de Polanco al AICM un viernes a las 6 PM
si hay marcha?''} y el agente razona sobre la petición completa.
\item \textbf{Voz (manos libres):} pulsa el micrófono, habla la
consulta y recibe la respuesta sintetizada en
\textless 30 palabras, comprensible en una sola escucha mientras
conduce.
\end{enumerate}

\subsection{Para el operador de flota}

El operador accede a \VialAI\ como herramienta de despacho:
puede ingresar múltiples rutas del día, activar perturbaciones
contextuales vigentes (p.~ej., ``marcha 9 de marzo'', ``cierre Línea~12
del Metro'') y obtener el IC de cada ruta para priorizar el despacho
de acuerdo con la confiabilidad real, no con el ETA nominal.

% ═══════════════════════ §3 EJEMPLO CONDUCTOR ══════════════════
\section{Ejemplo de funcionamiento: conductor ``Miguel''}

Se toma como caso de uso prototípico a Miguel, \num{28}~años,
repartidor para Uber Eats y Rappi, \num{22} entregas diarias en
promedio, smartphone Android de gama media. Necesita decidir si acepta
un pedido con origen en \textbf{Polanco} y entrega en
\textbf{Santa Fe} para las 08:15 de un martes hábil.

\subsubsection{Paso 1: consulta}

Miguel escribe en el chat de \VialAI:

\begin{quote}
\itshape ``¿Cuánto tardo de Polanco a Santa Fe? Necesito llegar antes
de las 9. Son las 8:15.''
\end{quote}

\subsubsection{Paso 2: procesamiento del motor estocástico}

\VialAI\ llama internamente a \texttt{predecir\_tiempo} con los
parámetros inferidos. El motor Monte Carlo ejecuta \num{10000}
trayectorias sobre la cadena de Markov calibrada con datos del C5
\CDMX\ (2018--2024) y las velocidades actuales de TomTom Traffic
Stats. No hay perturbación activa, por lo que la matriz de transición
base (día hábil, pico AM) se aplica sin modificar.

\subsubsection{Paso 3: resultado numérico}

\begin{table}[H]
\centering
\small
\renewcommand{\arraystretch}{1.4}
\begin{tabular}{@{}lrll@{}}
\toprule
\rowcolor{vialNavy!10}
\textbf{Métrica} & \textbf{Valor} & \textbf{Interpretación} & \textbf{Semáforo} \\
\midrule
$P_{10}$ (escenario optimista) & 19 min & Llegas en 19 min si el tráfico fluye bien & --- \\
$P_{50}$ (mediana esperada)    & 27 min & Lo más probable, llegada a las 08:42 & --- \\
$P_{90}$ (escenario adverso)   & 44 min & Si hay congestión fuerte, llegas a las 08:59 & --- \\
\midrule
\rowcolor{vialAmber!15}
Índice de Confiabilidad        & 0.93 & Alta variabilidad en esta ruta/hora & \textbf{Amarillo} \\
\bottomrule
\end{tabular}
\caption{Resultado del motor Monte Carlo para la ruta Polanco$\to$Santa Fe,
         martes, 08:15 AM. IC~$= (P_{90} - P_{10}) / P_{50} = (44-19)/27 \approx 0.93$.}
\label{tab:ejemplo_miguel}
\end{table}

\subsubsection{Paso 4: recomendación del agente}

\VialAI\ razona sobre la distribución y responde en lenguaje natural:

\begin{casebox}[Respuesta de \VialAI\ a Miguel]
\itshape
``Ruta de 17.3~km. El tiempo más probable es 27~minutos, pero hay
variabilidad alta (IC = amarillo): en el 10\% de los casos tardas
hasta 44 minutos. \textbf{Si necesitas estar en Santa Fe antes de las
9, te recomiendo salir ya}: tienes margen ajustado. Si el tráfico está
fluido, llegas a las 8:42. Si hay congestionamiento, podrías llegar
justo a las 8:59. No esperes.''
\end{casebox}

Miguel tiene un diagnóstico preciso. Puede tomar una decisión
informada, no una apuesta.

% ═══════════════════════ §4 EJEMPLO FLOTA ══════════════════════
\section{Ejemplo de funcionamiento: operador ``LogísticaMX''}

LogísticaMX opera una flota de \num{50} conductores en la \ZMVM.
Su gerente de operaciones activa la perturbación
\emph{``Marcha 9 de marzo''} antes del turno y solicita a \VialAI\
la evaluación de sus tres rutas críticas del día.

\begin{table}[H]
\centering
\small
\renewcommand{\arraystretch}{1.4}
\begin{tabular}{@{}lrrrrll@{}}
\toprule
\rowcolor{vialNavy!10}
\textbf{Ruta} & $P_{50}$ base & $P_{50}$ con perturbación
              & $P_{90}$ & IC & \textbf{Semáforo} & \textbf{Acción sugerida} \\
\midrule
Centro $\to$ Perisur & 38 min & 51 min & 79 min & 1.37 & \textcolor{red}{\textbf{Rojo}}
  & Posponer o ruta alterna \\
Polanco $\to$ AICM   & 29 min & 34 min & 49 min & 0.88 & \textcolor{vialAmber!80!black}{\textbf{Amarillo}}
  & Salir con \num{20} min extra \\
Ecatepec $\to$ Vallejo & 22 min & 23 min & 31 min & 0.36 & \textcolor{green!60!black}{\textbf{Verde}}
  & Despachar con normalidad \\
\bottomrule
\end{tabular}
\caption{Panel de rutas de LogísticaMX con perturbación activa ``Marcha 9 de marzo''.
         La ruta Centro$\to$Perisur cruza la zona de conflicto; la tercera ruta no
         se ve afectada.}
\label{tab:flota}
\end{table}

Con esta información, el gerente pospone el despacho de la ruta roja
y asigna buffers de tiempo extra a la amarilla. Ninguna decisión se
toma a ciegas.

% ═══════════════════════ §5 SEGMENTACIÓN ═══════════════════════
\section{Segmentación de mercado y modelo de negocio}

\VialAI\ segmenta el mercado en tres niveles con características de
adopción y disposición de pago distintas:

\begin{table}[H]
\centering
\small
\renewcommand{\arraystretch}{1.5}
\begin{tabularx}{\linewidth}{@{}l l X r r@{}}
\toprule
\rowcolor{vialNavy!10}
\textbf{Segmento} & \textbf{Perfil} & \textbf{Propuesta} & \textbf{Universo \ZMVM} & \textbf{Precio} \\
\midrule
Tier~1 (base) & Conductor individual & Predicción + agente de voz gratis
  & \num{180000} conductores & Gratuito \\
Tier~1 Plus & Conductor individual & Todo lo anterior + voz premium,
  sin anuncios, historial ampliado
  & Subconjunto de Tier~1 & \$49 MXN/mes \\
Tier~2 (B2B) & Operador de flota & Panel de flota, IC por ruta,
  alertas de perturbación, API de integración
  & \num{650} operadores formales & \$199 MXN/mes/conductor \\
Tier~3 (enterprise) & Empresa con SLA contractual & Calibración
  dedicada, SLA 99.5\%, soporte directo
  & Decenas de empresas & Cotización personalizada \\
\bottomrule
\end{tabularx}
\caption{Segmentación y modelo de precios de \VialAI. El Tier~1 gratuito
         sirve como \emph{wedge} de adopción masiva; el Tier~2 es el motor
         de ingresos principal.}
\label{tab:tiers}
\end{table}

La estrategia de adopción es \textbf{freemium con ancla masiva}:
se construye primero la base de Tier~1 (conductores individuales),
se genera telemetría de uso consentida para refinar el modelo, y con
el modelo validado se oferta el Tier~2 a operadores de flota.
El punto de equilibrio operativo estimado es de
\textbf{\num{800} usuarios pagos combinados entre Tier~1 Plus y Tier~2}.

% ═══════════════════════ §6 IMPACTO ECONÓMICO ══════════════════
\section{Cuantificación del impacto económico}

\subsection{Para una PyME de última milla típica (50 conductores)}

Las empresas de entrega de última milla incurren en costos ocultos
por demoras que no están en el presupuesto visible: bonos de SLA
pagados al cliente por incumplimiento, re-entregas, sobretiempo del
conductor y rotación debida a frustración operativa. Estimamos esos
costos para una PyME con \num{50} conductores activos en la \ZMVM:

\begin{table}[H]
\centering
\small
\renewcommand{\arraystretch}{1.5}
\begin{tabular}{@{}lrr@{}}
\toprule
\rowcolor{vialNavy!10}
\textbf{Concepto} & \textbf{Costo estimado / año} & \textbf{Supuesto} \\
\midrule
Penalizaciones SLA (incumplimientos) & \$210\,000 MXN & 3\% de entregas fuera de ventana \\
Re-entregas por conductor ausente & \$180\,000 MXN & 8 re-entregas/semana × \$430 costo \\
Sobretiempo no previsto & \$260\,000 MXN & 2 hr/conductor/semana × \$250/hr \\
\midrule
\textbf{Sobrecosto total estimado} & \textbf{\$650\,000 MXN} & --- \\
\midrule\midrule
Suscripción \VialAI\ Tier~2 (50 conductores) & \$119\,400 MXN & \$199 × 50 × 12 \\
\textbf{Ahorro neto anual} & \textbf{\$530\,600 MXN} & ROI $\approx$ 445\% \\
\bottomrule
\end{tabular}
\caption{Cálculo de impacto económico para una PyME de última milla con 50 conductores.
         Cifras basadas en tarifas de SLA de Uber Eats / Rappi 2024 y encuestas
         del Observatorio Laboral 2024.}
\label{tab:roi}
\end{table}

% ═══════════════════════ §7 COMPARATIVA ════════════════════════
\section{Comparativa competitiva}

\begin{table}[H]
\centering
\small
\renewcommand{\arraystretch}{1.5}
\begin{tabularx}{\linewidth}{@{}X c c c c@{}}
\toprule
\rowcolor{vialNavy!10}
\textbf{Capacidad} & \textbf{\VialAI} & \textbf{Google Maps}
                   & \textbf{Waze} & \textbf{TomTom} \\
\midrule
Estimación puntual (ETA) & \checkmark & \checkmark & \checkmark & \checkmark \\
Bandas de incertidumbre ($P_{10}$/$P_{90}$) & \textbf{\checkmark} & $\times$ & $\times$ & $\times$ \\
Índice de Confiabilidad & \textbf{\checkmark} & $\times$ & $\times$ & $\times$ \\
Agente conversacional contextual & \textbf{\checkmark} & $\times$ & $\times$ & $\times$ \\
Entrada/salida por voz (manos libres) & \textbf{\checkmark} & Parcial & Parcial & $\times$ \\
Perturbaciones locales \ZMVM\ & \textbf{\checkmark} & $\times$ & Parcial & $\times$ \\
MAPE validado en rutas \ZMVM & \textbf{11.9\%} & n/d & n/d & 19.0\% \\
Precio para conductores & Gratis & Gratis & Gratis & Pago \\
API de integración para flotas & \textbf{\checkmark} & \checkmark & $\times$ & \checkmark \\
\bottomrule
\end{tabularx}
\caption{Comparativa funcional de \VialAI\ contra competidores directos e indirectos.
         \VialAI\ es el único sistema que cuantifica la incertidumbre del viaje y
         además supera al líder comercial (TomTom) en precisión de estimación central.}
\label{tab:comparativa}
\end{table}

La ventaja competitiva de \VialAI\ no reside en una sola función:
reside en la \textbf{combinación} de rigor probabilístico (algo que
ningún competidor ofrece hoy), interfaz de voz específicamente
diseñada para el contexto del conductor en movimiento, y calibración
con fuentes de datos locales de la \ZMVM\ que los competidores
globales no utilizan.

% ═══════════════════════ §8 SIGUIENTES PASOS ═══════════════════
\section{Siguientes pasos}

El roadmap de los seis meses post-defensa se estructura en tres
hitos que escalan el sistema desde su estado actual de prototipo
validado hacia un producto desplegable:

\begin{enumerate}
\item \textbf{Mes 1--2: validación cuantitativa ampliada.}
Campaña de campo con $N \geq 500$ viajes etiquetados por exactitud,
reclutando \num{20}--\num{30} conductores profesionales del Tier~1.
Recalibración de la matriz de Markov con observaciones reales y
publicación de resultados con intervalos de confianza estrechos.

\item \textbf{Mes 3--4: versión móvil nativa.}
Migración de la interfaz web Streamlit a una aplicación Android
(Flutter o Kivy) apta para montarse en manillar. Integración de GPS
continuo para eliminar la necesidad de pulsar ``Llegué'' manualmente.
Calibración de perfil específico para motocicleta.

\item \textbf{Mes 5--6: primer piloto B2B.}
Despliegue controlado con un operador del Tier~2 (flota de
\num{20}--\num{50} conductores). Métricas de éxito: reducción de
demoras fuera de SLA $\geq 15\%$, satisfacción del conductor
$\geq 4/5$ estrellas, costo por entrega estable o decreciente.
Formalización del contrato de suscripción y estructura de facturación.
\end{enumerate}

\begin{tipbox}[Por qué este nicho y no otro]
El conductor de última milla de la \ZMVM\ cumple simultáneamente tres
condiciones que ningún otro segmento reúne: (i) el dolor de la
incertidumbre es \emph{cuantificable en pesos} (bonos de SLA,
re-entregas), (ii) la interfaz de voz es una \emph{restricción real de
seguridad}, no un lujo, y (iii) las fuentes abiertas del C5 y TomTom
cubren bien las zonas donde operan. Ningún otro nicho cumple las tres
condiciones al mismo tiempo.
\end{tipbox}

\end{document}
